\documentclass[12pt]{article}

\usepackage[utf8]{inputenc}
\usepackage[T1]{fontenc}
\usepackage{lmodern}
\usepackage[a4paper,margin=1in]{geometry}
\usepackage{xcolor}
\usepackage{listings}
\usepackage{setspace}
\usepackage{enumitem}
\usepackage{hyperref}
\usepackage{fancyhdr}
\usepackage{titlesec}

\definecolor{ash}{RGB}{235,235,235}

\hypersetup{
    colorlinks=false
}

\lstdefinestyle{codeStyle}{
    backgroundcolor=\color{ash},
    basicstyle=\ttfamily\small,
    keywordstyle=\bfseries,
    commentstyle=\itshape,
    stringstyle=,
    numbers=left,
    numberstyle=\scriptsize,
    numbersep=8pt,
    frame=single,
    breaklines=true,
    showstringspaces=false,
    tabsize=2,
    keepspaces=true
}

\lstset{style=codeStyle}

\pagestyle{fancy}
\fancyhf{}
\fancyhead[L]{Android Lab Report}
\fancyhead[R]{Humayra Adiba}
\fancyfoot[C]{\thepage}

\onehalfspacing

\begin{document}

\begin{titlepage}
\begin{center}
    \vspace*{1cm}
    {\Huge \textbf{Mobile App Development Report}}\\[0.4cm]
    {\Large Java + Android Studio + Wireless Debugging}\\[1.0cm]

    \rule{0.9\textwidth}{0.8pt}\\[0.6cm]
    \textbf{Name:} Humayra Adiba\\
    \textbf{Student ID:} 22701050\\
    \textbf{Department:} Computer Science and Engineering\\
    \textbf{University:} University of Chittagong\\[0.6cm]
    \rule{0.9\textwidth}{0.8pt}\\[1.2cm]

    \textbf{Course:} Mobile App Development Lab\\
    \textbf{Platform:} Windows + Android Studio + Java\\
    \textbf{Device Testing:} Real Android Phone (Wireless Debugging)\\[1.5cm]

    \today
\end{center}
\end{titlepage}

\tableofcontents
\newpage

\section*{Abstract}
\addcontentsline{toc}{section}{Abstract}
This report documents Android app development using Java on Windows. Android Studio and Android SDK were installed without emulator tools. A physical Android phone was used for build and test through wireless debugging. For every problem, one XML layout file and one Java activity file were implemented and then executed on the phone. The final outputs were verified on the real device.

\section{Development Environment}
\subsection{System and Tools}
\begin{itemize}[leftmargin=*]
    \item Operating System: Windows
    \item IDE: Android Studio
    \item Language: Java
    \item SDK: Android SDK (without emulator package)
    \item Testing Device: Personal Android phone
\end{itemize}

\subsection{Why No Emulator}
The emulator package was not installed. Therefore, all builds and run operations were tested directly on the phone. This provided real-device performance and avoided emulator setup overhead.

\section{Wireless Debugging Setup}
\subsection{Phone Configuration}
\begin{enumerate}[leftmargin=*]
    \item Enable \textbf{Developer options}.
    \item Enable \textbf{USB debugging}.
    \item Enable \textbf{Wireless debugging}.
    \item Connect phone and laptop to the same Wi-Fi network.
\end{enumerate}

\subsection{ADB Pair and Connect Commands}
\begin{lstlisting}[language=bash,caption={Wireless debugging commands}]
adb pair <PHONE_IP>:<PAIR_PORT>
# Enter pairing code shown on phone

adb connect <PHONE_IP>:<DEBUG_PORT>
adb devices
\end{lstlisting}

If the device appears in \texttt{adb devices}, Android Studio can run apps directly on the phone using Wi-Fi.

\section{Problem List}
\begin{enumerate}[leftmargin=*]
    \item Hello Adiba
    \item Temperature Converter
    \item BMI Converter
    \item NameApp
    \item Square App
\end{enumerate}

\section{Problem 1: Hello Adiba}
\subsection{Goal}
Display a greeting message for Adiba.

\subsection{XML Layout (activity\_main.xml)}
\begin{lstlisting}[language=XML]
<?xml version="1.0" encoding="utf-8"?>
<LinearLayout xmlns:android="http://schemas.android.com/apk/res/android"
    android:layout_width="match_parent"
    android:layout_height="match_parent"
    android:gravity="center"
    android:orientation="vertical"
    android:padding="24dp">

    <TextView
        android:id="@+id/tvHello"
        android:layout_width="wrap_content"
        android:layout_height="wrap_content"
        android:text="Hello Adiba"
        android:textSize="30sp"
        android:textStyle="bold" />

</LinearLayout>
\end{lstlisting}

\subsection{Java Code (MainActivity.java)}
\begin{lstlisting}[language=Java]
package com.example.helloadiba;

import androidx.appcompat.app.AppCompatActivity;
import android.os.Bundle;

public class MainActivity extends AppCompatActivity {
    @Override
    protected void onCreate(Bundle savedInstanceState) {
        super.onCreate(savedInstanceState);
        setContentView(R.layout.activity_main);
    }
}
\end{lstlisting}

\subsection{Output}
The phone screen shows the text: \textbf{Hello Adiba}.

\section{Problem 2: Temperature Converter}
\subsection{Goal}
Convert Celsius to Fahrenheit.

\subsection{XML Layout (activity\_main.xml)}
\begin{lstlisting}[language=XML]
<?xml version="1.0" encoding="utf-8"?>
<LinearLayout xmlns:android="http://schemas.android.com/apk/res/android"
    android:layout_width="match_parent"
    android:layout_height="match_parent"
    android:orientation="vertical"
    android:padding="20dp">

    <EditText
        android:id="@+id/etCelsius"
        android:layout_width="match_parent"
        android:layout_height="wrap_content"
        android:hint="Enter Celsius"
        android:inputType="numberDecimal" />

    <Button
        android:id="@+id/btnConvert"
        android:layout_width="match_parent"
        android:layout_height="wrap_content"
        android:layout_marginTop="12dp"
        android:text="Convert" />

    <TextView
        android:id="@+id/tvResult"
        android:layout_width="wrap_content"
        android:layout_height="wrap_content"
        android:layout_marginTop="14dp"
        android:text="Result will appear here"
        android:textSize="18sp" />

</LinearLayout>
\end{lstlisting}

\subsection{Java Code (MainActivity.java)}
\begin{lstlisting}[language=Java]
package com.example.temperatureconverter;

import androidx.appcompat.app.AppCompatActivity;
import android.os.Bundle;
import android.widget.Button;
import android.widget.EditText;
import android.widget.TextView;

public class MainActivity extends AppCompatActivity {
    @Override
    protected void onCreate(Bundle savedInstanceState) {
        super.onCreate(savedInstanceState);
        setContentView(R.layout.activity_main);

        EditText etCelsius = findViewById(R.id.etCelsius);
        Button btnConvert = findViewById(R.id.btnConvert);
        TextView tvResult = findViewById(R.id.tvResult);

        btnConvert.setOnClickListener(v -> {
            String input = etCelsius.getText().toString().trim();
            if (input.isEmpty()) {
                tvResult.setText("Please enter Celsius value");
                return;
            }

            double celsius = Double.parseDouble(input);
            double fahrenheit = (celsius * 9.0 / 5.0) + 32.0;
            tvResult.setText(String.format("Fahrenheit: %.2f", fahrenheit));
        });
    }
}
\end{lstlisting}

\subsection{Output}
After pressing Convert, the equivalent Fahrenheit value appears on screen.

\section{Problem 3: BMI Converter}
\subsection{Goal}
Calculate BMI from weight (kg) and height (m), then show BMI category.

\subsection{XML Layout (activity\_main.xml)}
\begin{lstlisting}[language=XML]
<?xml version="1.0" encoding="utf-8"?>
<LinearLayout xmlns:android="http://schemas.android.com/apk/res/android"
    android:layout_width="match_parent"
    android:layout_height="match_parent"
    android:orientation="vertical"
    android:padding="20dp">

    <EditText
        android:id="@+id/etWeight"
        android:layout_width="match_parent"
        android:layout_height="wrap_content"
        android:hint="Weight (kg)"
        android:inputType="numberDecimal" />

    <EditText
        android:id="@+id/etHeight"
        android:layout_width="match_parent"
        android:layout_height="wrap_content"
        android:layout_marginTop="10dp"
        android:hint="Height (m)"
        android:inputType="numberDecimal" />

    <Button
        android:id="@+id/btnBmi"
        android:layout_width="match_parent"
        android:layout_height="wrap_content"
        android:layout_marginTop="12dp"
        android:text="Calculate BMI" />

    <TextView
        android:id="@+id/tvBmiResult"
        android:layout_width="wrap_content"
        android:layout_height="wrap_content"
        android:layout_marginTop="14dp"
        android:text="BMI result"
        android:textSize="18sp" />

</LinearLayout>
\end{lstlisting}

\subsection{Java Code (MainActivity.java)}
\begin{lstlisting}[language=Java]
package com.example.bmiconverter;

import androidx.appcompat.app.AppCompatActivity;
import android.os.Bundle;
import android.widget.Button;
import android.widget.EditText;
import android.widget.TextView;

public class MainActivity extends AppCompatActivity {
    @Override
    protected void onCreate(Bundle savedInstanceState) {
        super.onCreate(savedInstanceState);
        setContentView(R.layout.activity_main);

        EditText etWeight = findViewById(R.id.etWeight);
        EditText etHeight = findViewById(R.id.etHeight);
        Button btnBmi = findViewById(R.id.btnBmi);
        TextView tvBmiResult = findViewById(R.id.tvBmiResult);

        btnBmi.setOnClickListener(v -> {
            String wText = etWeight.getText().toString().trim();
            String hText = etHeight.getText().toString().trim();

            if (wText.isEmpty() || hText.isEmpty()) {
                tvBmiResult.setText("Please enter weight and height");
                return;
            }

            double weight = Double.parseDouble(wText);
            double height = Double.parseDouble(hText);
            if (height <= 0) {
                tvBmiResult.setText("Height must be greater than zero");
                return;
            }

            double bmi = weight / (height * height);
            String status;
            if (bmi < 18.5) {
                status = "Underweight";
            } else if (bmi < 25.0) {
                status = "Normal";
            } else if (bmi < 30.0) {
                status = "Overweight";
            } else {
                status = "Obese";
            }

            tvBmiResult.setText(String.format("BMI: %.2f (%s)", bmi, status));
        });
    }
}
\end{lstlisting}

\subsection{Output}
The app shows BMI value and category after calculation.

\section{Problem 4: NameApp}
\subsection{Goal}
Take a name as input and display a personalized message.

\subsection{XML Layout (activity\_main.xml)}
\begin{lstlisting}[language=XML]
<?xml version="1.0" encoding="utf-8"?>
<LinearLayout xmlns:android="http://schemas.android.com/apk/res/android"
    android:layout_width="match_parent"
    android:layout_height="match_parent"
    android:orientation="vertical"
    android:padding="20dp">

    <EditText
        android:id="@+id/etName"
        android:layout_width="match_parent"
        android:layout_height="wrap_content"
        android:hint="Enter your name" />

    <Button
        android:id="@+id/btnShow"
        android:layout_width="match_parent"
        android:layout_height="wrap_content"
        android:layout_marginTop="12dp"
        android:text="Show Message" />

    <TextView
        android:id="@+id/tvNameResult"
        android:layout_width="wrap_content"
        android:layout_height="wrap_content"
        android:layout_marginTop="14dp"
        android:textSize="18sp" />

</LinearLayout>
\end{lstlisting}

\subsection{Java Code (MainActivity.java)}
\begin{lstlisting}[language=Java]
package com.example.nameapp;

import androidx.appcompat.app.AppCompatActivity;
import android.os.Bundle;
import android.widget.Button;
import android.widget.EditText;
import android.widget.TextView;

public class MainActivity extends AppCompatActivity {
    @Override
    protected void onCreate(Bundle savedInstanceState) {
        super.onCreate(savedInstanceState);
        setContentView(R.layout.activity_main);

        EditText etName = findViewById(R.id.etName);
        Button btnShow = findViewById(R.id.btnShow);
        TextView tvNameResult = findViewById(R.id.tvNameResult);

        btnShow.setOnClickListener(v -> {
            String name = etName.getText().toString().trim();
            if (name.isEmpty()) {
                tvNameResult.setText("Please enter a name");
            } else {
                tvNameResult.setText("Welcome, " + name + "!");
            }
        });
    }
}
\end{lstlisting}

\subsection{Output}
The app displays a custom welcome message using the entered name.

\section{Problem 5: Square App}
\subsection{Goal}
Input a number and show its square.

\subsection{XML Layout (activity\_main.xml)}
\begin{lstlisting}[language=XML]
<?xml version="1.0" encoding="utf-8"?>
<LinearLayout xmlns:android="http://schemas.android.com/apk/res/android"
    android:layout_width="match_parent"
    android:layout_height="match_parent"
    android:orientation="vertical"
    android:padding="20dp">

    <EditText
        android:id="@+id/etNumber"
        android:layout_width="match_parent"
        android:layout_height="wrap_content"
        android:hint="Enter a number"
        android:inputType="numberDecimal" />

    <Button
        android:id="@+id/btnSquare"
        android:layout_width="match_parent"
        android:layout_height="wrap_content"
        android:layout_marginTop="12dp"
        android:text="Calculate Square" />

    <TextView
        android:id="@+id/tvSquareResult"
        android:layout_width="wrap_content"
        android:layout_height="wrap_content"
        android:layout_marginTop="14dp"
        android:textSize="18sp" />

</LinearLayout>
\end{lstlisting}

\subsection{Java Code (MainActivity.java)}
\begin{lstlisting}[language=Java]
package com.example.squareapp;

import androidx.appcompat.app.AppCompatActivity;
import android.os.Bundle;
import android.widget.Button;
import android.widget.EditText;
import android.widget.TextView;

public class MainActivity extends AppCompatActivity {
    @Override
    protected void onCreate(Bundle savedInstanceState) {
        super.onCreate(savedInstanceState);
        setContentView(R.layout.activity_main);

        EditText etNumber = findViewById(R.id.etNumber);
        Button btnSquare = findViewById(R.id.btnSquare);
        TextView tvSquareResult = findViewById(R.id.tvSquareResult);

        btnSquare.setOnClickListener(v -> {
            String input = etNumber.getText().toString().trim();
            if (input.isEmpty()) {
                tvSquareResult.setText("Please enter a number");
                return;
            }

            double n = Double.parseDouble(input);
            double square = n * n;
            tvSquareResult.setText(String.format("Square: %.2f", square));
        });
    }
}
\end{lstlisting}

\subsection{Output}
The square result is shown immediately after pressing the button.

\section{Testing and Results on Phone}
All five apps were run on a real Android phone using wireless debugging. Each app compiled successfully and produced expected output on the device screen.

\section{Conclusion}
This work demonstrates complete Android app development using Java and XML on Windows without an emulator. Wireless debugging made it possible to build, deploy, and test directly on the phone. The practical approach of writing one XML file and one Java file for each problem helped in understanding Android UI design, event handling, and basic computational logic.

\end{document}
